\documentclass[11pt,a4paper]{article}
\usepackage{webclases-lesson}
\lessonSetup{T3}{v0.2}{T3 - Leyes de Newton}

\begin{document}
\shorthandoff{<>}
\color{Ink}

\coverHeader{T3 --- Leyes de Newton}{Interacciones, diagramas de cuerpo libre y ecuaciones de movimiento.}

\begin{lessoncard}{Prop\'osito}
La din\'amica pregunta por qu\'e cambia el movimiento. En esta lecci\'on aislamos cuerpos, nombramos interacciones y traducimos la resultante a aceleraci\'on. Usaremos $g=10\mpss$ como aproximaci\'on did\'actica.
\begin{keybox}
\textbf{Clave de lectura.} La resultante no es una fuerza nueva: es la suma vectorial de las fuerzas reales que otros cuerpos ejercen sobre el objeto elegido.
\end{keybox}
\end{lessoncard}

\begin{lessonindex}
\indexentry{sec:b1}{1. Reconocer fuerzas y sus efectos}
\indexentry{sec:b2}{2. Resolver con diagramas y ejes}
\indexentry{sec:b3}{3. Fuerzas que hacen girar}
\end{lessonindex}
\precisionnote

\block{1}{Reconocer fuerzas y sus efectos}{sec:b1}

\begin{lessoncard}{Fuerzas y cuerpo libre}
Una fuerza representa una interacci\'on de un agente sobre un receptor: Tierra sobre bloque, mesa sobre bloque, cuerda sobre bloque. Al aislar un cuerpo, se dibujan solo las fuerzas que act\'uan sobre \'el. Velocidad y aceleraci\'on pueden aparecer en un esquema cinem\'atico separado, pero no son fuerzas.

\begin{center}
\begin{tikzpicture}[x=1cm,y=1cm]
  \draw[fill=Panel,draw=Border] (-.2,0) rectangle (3.8,.18);
  \draw[fill=BlueSoft,draw=BrandBlue] (1.2,.18) rectangle (2.4,1.0);
  \node at (1.8,.58) {\small bloque};
  \draw[force] (1.8,1.0) -- (1.8,1.9) node[above] {\small $\vec N$ mesa sobre bloque};
  \draw[force] (1.8,.18) -- (1.8,-.72) node[below] {\small $\vec P$ Tierra sobre bloque};
  \draw[force] (2.4,.6) -- (3.25,.6) node[right] {\small $\vec F$ mano};
  \draw[guide] (5.0,-.7) rectangle (7.4,1.9);
  \node[text=Muted] at (6.2,2.5) {\small cuerpo aislado};
  \draw[fill=BlueSoft,draw=BrandBlue] (5.8,.35) rectangle (6.6,1.05);
  \draw[force] (6.2,1.05) -- (6.2,1.75) node[above] {\small $\vec N$};
  \draw[force] (6.2,.35) -- (6.2,-.35) node[below] {\small $\vec P$};
  \draw[force] (6.6,.7) -- (7.25,.7) node[right] {\small $\vec F$};
\end{tikzpicture}
\end{center}
\end{lessoncard}

\begin{lessoncard}{Resultante y leyes}
En un sistema inercial, resultante cero significa velocidad vectorial constante, que puede ser nula. Para masa constante:
\begin{formulabox}
\[
\sum \vec F=m\vec a \qquad 1\N=1\kg\,\mathrm{m/s^2}
\]
\end{formulabox}
Una fuerza aislada no determina por s\'i sola la aceleraci\'on si otras la compensan. La primera ley sirve para reconocer el marco inercial; la segunda cuantifica el cambio cuando la resultante no es cero.

\heading{Acci\'on y reacci\'on: dos cuerpos distintos}
Si A ejerce una fuerza sobre B, B ejerce otra de igual m\'odulo y sentido opuesto sobre A. Act\'uan sobre cuerpos distintos, por eso no se cancelan dentro del diagrama de un solo cuerpo. Peso y normal del mismo bloque no son un par de acci\'on-reacci\'on.
\end{lessoncard}

\heading{Qu\'e significa un sistema inercial}
Un marco inercial es aquel en el que un cuerpo libre de resultante mantiene
velocidad constante. Para los problemas del aula tomamos aproximadamente el suelo
como inercial. Un autob\'us que frena no lo es: respecto a \'el un objeto puede
cambiar su movimiento aunque no aparezca la interacci\'on que uno esperaba.
No a\~nadimos aqu\'i fuerzas ficticias; elegimos el marco del suelo.

\heading{Masa, peso y normal}
La masa mide inercia y se expresa en kg. El peso es la fuerza gravitatoria
$\vec P=m\vec g$, en N. La normal aparece si hay contacto y empuja perpendicularmente
a la superficie. Su valor sale de Newton; una superficie no puede tirar del objeto.
Si un c\'alculo da $N<0$, hay que revisar la hip\'otesis de contacto.

\begin{examplebox}{mantener una velocidad ya existente}
\origin{Adaptado del original T3, p.\ 5.}
\step{Planteamiento} Un bloque de $0,250\kg$ se mueve sobre mesa lisa. Act\'uan $F_1=3,00\N$ a la derecha y $F_2=1,00\N$ a la izquierda. Halla su aceleraci\'on y qu\'e $F_1$ mantendr\'ia su velocidad.
\solutionblock

\step{Predicci\'on} Como la resultante apunta a la derecha, la aceleraci\'on tambi\'en apunta a la derecha.

\step{Modelo} Elegimos $+x$ a la derecha y aplicamos $\sum F_x=ma_x$.
\step{C\'alculo}
\[
\sum F_x=3,00-1,00=2,00\N,\qquad a_x=\frac{2,00}{0,250}=8,00\mpss.
\]
Para mantener una velocidad ya existente de $+1,00\mps$ hace falta $\sum F_x=0$, luego $F_1=1,00\N$.

\step{Interpretaci\'on} Esa fuerza no crea instant\'aneamente una velocidad de $1\mps$ desde reposo: solo mantiene la velocidad que ya ten\'ia. Si inicialmente va a la izquierda, una resultante positiva primero reduce su rapidez.
\end{examplebox}

\begin{exercisebox}{Corrige el diagrama}
En un bloque en reposo sobre una mesa se dibujan peso, normal y una tercera fuerza llamada resultante. Adem\'as se afirma que peso y normal son acci\'on y reacci\'on. Corrige las dos afirmaciones y nombra el receptor de cada reacci\'on, con un diagrama y dos frases.
\solutionblock La resultante es cero y no se a\~nade como fuerza real. La reacci\'on al peso es la fuerza gravitatoria del bloque sobre la Tierra; la reacci\'on a la normal es la fuerza del bloque sobre la mesa.
\end{exercisebox}


\block{2}{Resolver con diagramas y ejes}{sec:b2}

\heading{Un m\'etodo que se puede repetir}
Primero elige un cuerpo y enumera sus interacciones; despu\'es dibuja fuerzas
desde ese cuerpo. Proyecta sobre ejes convenientes y escribe
$\sum F_x=ma_x$, $\sum F_y=ma_y$. Una direcci\'on sin aceleraci\'on no implica que
no haya fuerzas, sino que sus componentes se compensan.
Las componentes del peso sustituyen al vector al proyectar: no se suman otra vez
junto al peso completo.

\begin{examplebox}{la normal no siempre es el peso}
\origin{Elaboraci\'on propia.}
\step{Planteamiento} Un bloque de $2,00\kg$ est\'a sobre una mesa.
Se tira de \'el con $10,0\N$ a $30^\circ$ sobre la horizontal.
Calcula la normal si no hay aceleraci\'on vertical.
\solutionblock
\step{Modelo} En vertical act\'uan $N$, $F\sin30^\circ$ y $-mg$.
La componente ascendente de la fuerza reduce el apoyo.
\step{C\'alculo} $N+10\sin30^\circ-20=0$, luego $N=15,0\N$.
\step{Interpretaci\'on} Con la misma fuerza hacia abajo a $30^\circ$, $N=25,0\N$.
Ni el peso ni la masa cambian. El contacto es compatible con $N>0$.
\end{examplebox}

\begin{lessoncard}{Componentes en un plano}
La normal es perpendicular a la superficie: no siempre vale $mg$. En un plano liso de \'angulo $\alpha$, conviene tomar $s$ paralelo al plano y $n$ perpendicular.
\[
N=mg\cos\alpha,\qquad \sum F_s=-mg\sin\alpha=ma_s.
\]
\begin{center}
\begin{tikzpicture}[x=1cm,y=1cm,font=\small]
\pgfmathsetmacro{\h}{5.4*tan(20)}
\draw[fill=Panel,draw=Border] (0,0)--(5.4,0)--(5.4,\h)--cycle;
\begin{scope}[rotate=20]
 \draw[fill=BlueSoft,draw=BrandBlue] (2.2,0) rectangle (3,.55);
 \coordinate (B) at (2.6,.275);
 \draw[force] (B)--++(0,{1.6*cos(20)}) node[above] {$N$};
 \draw[guide] (B)--++({-1.6*sin(20)},0) node[left] {$mg\sin20^\circ$};
 \draw[guide] (B)--++(0,{-1.6*cos(20)}) node[right] {$mg\cos20^\circ$};
 \draw[axis] (3.4,.3)--(4.8,.3) node[right] {$+s$};
\end{scope}
\draw[force] (B)--++(0,-1.6) node[below] {$mg$};
\draw (.8,0) arc[start angle=0,end angle=20,radius=.8];
\node at (1.2,.18) {$20^\circ$};
\end{tikzpicture}
\end{center}
\end{lessoncard}

\begin{examplebox}{Subir y bajar el mismo plano}
\origin{Adaptado del original T3, pp.\ 6--7.}
\step{Planteamiento} Plano liso de $20^\circ$. Se lanza un bloque cuesta arriba con $v_0=2,60\mps$ desde $s=0$. Eje positivo hacia arriba del plano. Halla tiempo hasta parar, distancia sobre el plano y ganancia de altura.
\solutionblock

\step{Modelo} Solo act\'uan peso y normal; la aceleraci\'on paralela es $a_s=-g\sin20^\circ=-3,420\mpss$.
\[
v=v_0+a_st,\qquad s=v_0t+\frac12 a_st^2.
\]
\step{C\'alculo} En el punto m\'as alto $v=0$:
$t=-v_0/a_s$, $s=-v_0^2/(2a_s)$ y $h=s\sin20^\circ$. As\'i,
\[
t=0,760\units{s},\quad s=0,988\units{m},\quad h=s\sin20^\circ=0,338\units{m}.
\]
Al volver al origen, $v=-2,60\mps$ y el tiempo total es $1,52\units{s}$ si el plano contin\'ua.

\step{Interpretaci\'on} La aceleraci\'on apunta cuesta abajo tambi\'en mientras sube. En el punto superior $v=0$, pero $a$ no es cero.
\variation{Duplicar la masa no cambia $a$ ni $h$ en este modelo liso. En un descenso desde reposo por $30^\circ$ durante $0,600\units{m}$: $a=5,00\mpss$, $t=0,490\units{s}$ y $v=2,45\mps$.}
\end{examplebox}

\begin{lessoncard}{Cuerdas y poleas}
Una cuerda ideal tensa e inextensible y una polea ideal permiten usar la misma tensi\'on y relacionar aceleraciones. La tensi\'on siempre tira a lo largo de la cuerda.
Ideal significa cuerda de masa despreciable e inextensible, polea de inercia
despreciable y eje sin rozamiento. Que la cuerda no cambie de longitud impone
igual m\'odulo de aceleraci\'on en los dos tramos de este montaje.
No significa igual vector aceleraci\'on: uno es horizontal y otro vertical.
\end{lessoncard}

\begin{examplebox}{Dos masas, dos diagramas}
\origin{Adaptado del original T3, p.\ 10.}
\step{Planteamiento} Calcula aceleraci\'on, tensi\'on y normal para las siguientes masas, inicialmente en reposo.
$m_1=0,100\kg$ sobre mesa lisa y $m_2=0,200\kg$ colgante. Positivo: $m_1$ hacia la polea y $m_2$ hacia abajo.
\solutionblock
\begin{center}
\begin{tikzpicture}[x=1cm,y=1cm]
  \draw[fill=Panel,draw=Border] (-.4,0) rectangle (2.65,.15);
  \draw[fill=BlueSoft,draw=BrandBlue] (.6,.15) rectangle (1.35,.75);
  \draw[force] (1.35,.45) -- (2.35,.45) node[above] {\small $T$};
  \draw[force] (.98,.75) -- (.98,1.35) node[above] {\small $N$};
  \draw[force] (.98,.15) -- (.98,-.45) node[below] {\small $m_1g$};
  \draw (2.8,.15) circle (.30);
  \draw (1.35,.45) -- (2.8,.45) arc(90:0:.30) -- (3.1,-.65);
  \draw[fill=AmberSoft,draw=AmberLine] (2.8,-1.35) rectangle (3.4,-.65);
  \draw[force] (3.1,-.65) -- (3.1,-.15) node[right] {\small $T$};
  \draw[force] (3.1,-1.35) -- (3.1,-2.0) node[below] {\small $m_2g$};
\end{tikzpicture}
\end{center}
\step{Modelo} Cada masa tiene su propia ecuaci\'on. La cuerda vincula los
m\'odulos de aceleraci\'on; al sumar estas ecuaciones se elimina la inc\'ognita $T$.
\[
T=m_1a,\qquad m_2g-T=m_2a
\]
\step{C\'alculo}
\[
a=\frac{m_2g}{m_1+m_2}=6,67\mpss,\qquad T=m_1a=0,667\N.
\]
\step{Interpretaci\'on} La normal de la mesa sobre $m_1$ es $1,00\N$. $T<m_2g$ porque la fuerza neta sobre la masa colgante debe apuntar hacia abajo.
Si se impusiera $T=m_2g$, Newton dar\'ia $a=0$ para esa masa.
El modelo termina cuando una masa toca un tope o se pierde la tensi\'on.
\variation{Si ambas masas son $0,100\kg$, entonces $a=5,00\mpss$ y $T=0,500\N$. Las tensiones horizontal y vertical no son vectores que se cancelen directamente: se eliminan al sumar ecuaciones escalares compatibles.}
\end{examplebox}


\block{3}{Fuerzas que hacen girar}{sec:b3}

\begin{lessoncard}{Radial y tangencial}
En una trayectoria circular:
\begin{formulabox}
\[
\sum F_{\mathrm{radial}}=\frac{mv^2}{R}\quad\text{hacia el centro},\qquad
\sum F_{\mathrm{tangencial}}=m\frac{dv}{dt}.
\]
\end{formulabox}
``Centr\'ipeta'' describe la resultante radial, no una interacci\'on adicional. Si la rapidez es constante, la velocidad vectorial no lo es: cambia de direcci\'on.
\end{lessoncard}

\begin{examplebox}{P\'endulo c\'onico}
\origin{Geometr\'ia del original T3; datos propios.}
\step{Planteamiento} Esfera de $0,200\kg$ con cuerda de $1,00\units{m}$ que forma $30^\circ$ con la vertical. Giro horizontal uniforme sin resistencia del aire. Calcula radio, tensi\'on y rapidez.
\solutionblock
\step{Modelo} La altura no cambia: la componente vertical de $T$ compensa el peso.
Su componente horizontal apunta al centro. No a\~nadimos otra fuerza centr\'ipeta.
\begin{center}
\begin{tikzpicture}[x=1cm,y=1cm,font=\small]
\node at (.8,3.1) {vista lateral};
\coordinate (O) at (0,2.5); \coordinate (B) at (1.25,.33494);
\draw[guide] (O)--(0,0);
\draw[thick] (O)--(B);
\fill[BrandBlue] (B) circle (2.5pt);
\draw[force] (B)--++(120:1.5) node[above right] {$T$};
\draw[force] (B)--++(0,{-1.5*cos(30)}) node[below] {$mg$};
\draw[guide] (0,.33494)--(B) node[midway,below] {$R$};
\draw (0,1.8) arc[start angle=-90,end angle=-60,radius=.7];
\node at (.25,1.55) {$30^\circ$};
\begin{scope}[xshift=6.5cm,yshift=1cm]
 \node at (0,2.1) {vista superior};
 \draw[BrandBlue] (0,0) circle (1.25);
 \fill[BrandBlue] (1.25,0) circle (2.5pt);
 \draw[vector] (1.25,0)--(1.25,1.4) node[above] {$\vec v$};
 \draw[vectorgreen] (1.25,0)--(.2,0) node[below] {$\vec a_n$};
 \draw[guide] (0,0)--(0,-1.25) node[midway,left] {$R$};
\end{scope}
\end{tikzpicture}
\end{center}
\[
T\cos\theta=mg,\qquad T\sin\theta=\frac{mv^2}{R}.
\]
\step{C\'alculo} $R=L\sin30^\circ=0,500\units{m}$; al dividir las dos
ecuaciones, $\tan\theta=v^2/(Rg)$.
\[
T=\frac{mg}{\cos30^\circ}=2,31\N,\qquad
v=\sqrt{Rg\tan30^\circ}=1,70\mps.
\]
\step{Interpretaci\'on} El radio de la circunferencia no es la longitud de la cuerda: es su proyecci\'on horizontal.
\variation{Duplicar la masa duplica la tensi\'on, pero no cambia la rapidez para $L$ y $\theta$ fijos.}
\end{examplebox}

\begin{warnbox}
No confundas la normal de contacto $N$ con una ``componente normal'' de la aceleraci\'on. En din\'amica circular, primero identifica fuerzas reales y despu\'es decide qu\'e parte de su suma apunta hacia el centro.
\end{warnbox}

\begin{lessoncard}{Cierre}
Antes de aplicar una f\'ormula, escribe el cuerpo elegido, agentes de cada fuerza, ejes y condiciones. Si las fuerzas internas pertenecen al sistema completo, pueden desaparecer de la ecuaci\'on total; siguen existiendo en los diagramas individuales.
\end{lessoncard}

\begin{exercisebox}{qu\'e pasa al soltar la cuerda}
En el p\'endulo c\'onico anterior se rompe la cuerda. Dibuja la velocidad justo al
romperse y la aceleraci\'on inmediatamente despu\'es. Justifica en dos frases.
\solutionblock
La velocidad inicial es tangente al c\'irculo horizontal. Al desaparecer la
tensi\'on solo queda la gravedad: la aceleraci\'on es vertical hacia abajo.
El movimiento posterior es un tiro horizontal de T2, no una salida radial.
\end{exercisebox}

\heading{Fuentes y procedencia}
Los ejemplos de fuerzas opuestas, plano y polea adaptan el original T3;
el p\'endulo conserva su geometr\'ia con datos propios. Normal variable y
actividades conceptuales son de elaboraci\'on propia, igual que todos los diagramas.
\source{https://fisquiweb.es/Apuntes/apun1B.htm}{FisQuiWeb / IES La Magdalena: inspiraci\'on principal, 1.\textordmasculine\ Bachillerato.}
Contraste externo: interpretar la fuerza centr\'ipeta como resultante radial,
identificando la interacci\'on que la proporciona.
\source{https://openstax.org/books/university-physics-volume-1/pages/6-3-centripetal-force}{OpenStax, University Physics 1, 6.3: Centripetal Force.}
\end{document}
